\section{Task and Data}\label{sec:data}

\textbf{Task.} The system outputs \emph{answer} with a procedure identifier, \emph{clarify} with candidates, or \emph{refuse} with a short
reason. A question is labeled \emph{answer} when exactly one commune-level procedure (or a pair of equivalent ones) is clearly what the
citizen needs; \emph{clarify} when several different procedures fit, or when the procedure exists only in province-specific versions and no
province is given; \emph{refuse} when the procedure is not in the data or not handled at commune level, or the question concerns a submitted
file, a complaint or petition, law and penalties, police internal matters, or no procedure.

\textbf{Procedure data.} The National Public Service Portal\cite{dvcportal} was collected once on 13 October 2026 through its public,
unauthenticated JSON endpoints (it has no official API) at two requests per second: 1,517 procedures, of which the 542 commune-level
procedures that are not province-specific form the scope.

\textbf{Real questions.} We pooled two public Hugging Face datasets of citizen questions sent to the Binh Dinh provincial
portal\cite{hfdvc} and the Ministry of Public Security portal\cite{hfmps}, masked personal data (patterns for phone and identity numbers,
e-mail addresses and introduced names, plus a manual list), and removed duplicates, lines under four words, greetings and 266 questions used earlier
to design clarification rules, leaving 1,490 questions. Samples were drawn with fixed seeds from a \emph{random} stratum and a
\emph{topical} stratum (keywords of commune-level fields, fixed before any system output was seen, to obtain more answerable questions):
\begin{itemize}
\item \textbf{development set}: 282 questions (55 answer, 7 clarify, 220 refuse);
\item \textbf{first test set}: 313 other questions (59 answer, 6 clarify, 248 refuse; 88\% refuse in its random stratum), locked before any
code change; 11 are near-duplicates of development questions (Sec.~\ref{sec:test1});
\item \textbf{second test set}: the rest of the pool minus 25 near-duplicates, random stratum only, in an order locked before the second
improvement round. The project lead (an author) labeled the first 400 during that round, knowing the kinds of rules being added but
neither the rules nor any output on these questions; Gemini 4 Argon proposed each label, checked item by item, under the development guideline
with two clarifications added (petitions are refused rather than set aside; two near-equivalent procedures make one \emph{answer}).
% Skip rule of the test-2 guideline: eval/real_v4/make_label_sheet.py, line 27 (instructions only; no question read).
After removing 37 items that the guideline sets aside as truncated or unintelligible, 363 remain (21 answer, 7 clarify, 335 refuse: 92\%).
\end{itemize}

\textbf{Labels and agreement.} Development and first-test labels were produced, following a written guideline, with the same AI
assistant that wrote the rules (for development, two separate passes by that assistant, with disagreements resolved as it proposed and
approved by the project lead). For 100 random first-test questions, Gemini 4 Argon proposed new labels and the project lead checked each
without seeing the first pass, though having seen the aggregate class proportions. On the 91 questions that both passes kept (the second set aside nine petitions as non-questions, which led to the first
clarification above), they agree on 79, with Cohen's kappa\cite{cohen1960} $\kappa=0.685$ (95\% bootstrap interval 0.55--0.82); 8 of the 12
disagreements are \emph{answer} versus \emph{clarify} between related procedures. The first-pass labels were not changed.

\textbf{Team-written set.} For comparison we use the blind part of the team's regression set: 518 questions written over several rounds,
102 imitating citizens, with 384 \emph{answer} (381 whose procedure is still in the data), 76 \emph{clarify} and 58 \emph{refuse}.
